\section*{Hoofdstuk \ref{ch:b2:energy-geology} --- Energiegeologie}
% numbers from the Book 4 second-writer scratch calculator calc/ch30.py

\begin{solution}{exo:b2:energy-geology:1}
(a) 1.58: type~I; 0.95: type~III; 1.24: type~II; 1.12: gemengd type~II/III. Het kerogeen van type~III, met
0.95, levert vooral gas. (b) H/C $= 150/100 = 1.50$, O/C $= 8/100 = 0.08$. (c) Per \qty{100}{g}:
$77.4/12.011 = \qty{6.444}{mol}$ C, $9.8/1.008 = \qty{9.722}{mol}$ H, $8.2/15.999 =
\qty{0.513}{mol}$ O; H/C $= 1.51$ en O/C $= 0.0795$, in overeenstemming met de formule.
\end{solution}

\begin{solution}{exo:b2:energy-geology:2}
(a) 0.45: onrijp, onder het begin bij \qty{0.60}{\percent}; 0.68: levert olie; 1.1: nog in het
olievenster, dicht bij het einde ervan; 1.6: voorbij 1.20--\qty{1.30}{\percent}, dus olie valt uiteen tot methaan en het
gesteente levert gas; 2.5: alleen gas, waarbij de vorming doorgaat tot ongeveer \qty{4.8}{\percent}. (b) Nee: de
methaanvorming uit type~III begint rond \qty{0.73}{\percent}. (c) $F = (\ln R_o + 1.6)/3.7$: $(-0.511 +
1.6)/3.7 = 0.294$ en $(0.262 + 1.6)/3.7 = 0.503$.
\end{solution}

\begin{solution}{exo:b2:energy-geology:3}
(a) In het reservoir $6 \times 10^6 \times 25 \times 0.18 \times 0.70 = \qty{1.89e7}{m^3}$;
$\mathrm{STOIIP} = 1.89 \times 10^7/1.25 = \qty{1.51e7}{m^3}$. (b) $0.30 \times 1.512 \times 10^7 =
\qty{4.54e6}{m^3}$ en $0.40 \times 1.512 \times 10^7 = \qty{6.05e6}{m^3}$.
\end{solution}

\begin{solution}{exo:b2:energy-geology:4}
(a) $V = 10^7 \times 200 = \qty{2e9}{m^3}$ en $q_R = 2.5 \times 10^6 \times 2 \times 10^9 \times 60
= \qty{3.0e17}{J}$. (b) Onder \qty{90}{\celsius}: een \omterm{def:b2:energy-geology:geothermal}{lage-enthalpiebron}, voor rechtstreeks gebruik van de
warmte. (c) $P = 0.25 \times 3.0 \times 10^{17}/(30 \times \qty{3.156e7}{s}) = \qty{79.2}{MW}$ warmte.
\end{solution}

\begin{solution}{exo:b2:energy-geology:5}
(a) $k_2/k_1 = \exp[(E/R)(1/T_1 - 1/T_2)] = 2$ met $1/373.15 - 1/383.15 = \qty{6.994e-5}{K^{-1}}$
geeft $E = 8.314 \times \ln 2/(6.994 \times 10^{-5}) = \qty{82.4}{kJ/mol}$. (b) $\exp(26800 \times
6.994 \times 10^{-5}) = 6.52$. (c) De verdubbelingsregel, met een lagere effectieve energie, geeft
relatief meer gewicht aan de lange tijden bij lage temperatuur; voor een kerogeen van type~I met een hoge
activeringsenergie kondigt de TTI de olievorming te vroeg aan.
\end{solution}

\begin{solution}{exo:b2:energy-geology:6}
(a) $x = 1/2$ in de benadering geeft $(ART^2/\beta E)\,e^{-E/RT} = \ln 2$; logaritmen nemen geeft
$E/RT = \ln[AT^2/(\beta(E/R)\ln 2)]$, en de gegeven vorm volgt. Vanaf een willekeurige schatting rond
\qty{400}{K} en iterend met $\beta$ in kelvin per seconde: bij \qty{1}{\celsius\per\Myr}
($\beta = \qty{3.169e-14}{K/s}$) is $T = \qty{409.6}{K}$, \qty{136.4}{\celsius}; bij \qty{10}{K} per
minuut ($\beta = \qty{0.1667}{K/s}$, $\num{5.26e12}$ keer sneller) is $T = \qty{719.0}{K}$,
\qty{445.9}{\celsius}. (b) Het verschil is \qty{309.5}{K}; de schatting $6.259 \times
\ln(5.26 \times 10^{12}) = 6.259 \times 29.29 = \qty{183}{K}$ is veel te klein, omdat $RT^2/E$
zelf groeit met $T$ (tot \qty{19.3}{K} bij \qty{719}{K}): de lineaire verschuiving geldt alleen voor kleine
veranderingen van de snelheid.
\end{solution}

\begin{solution}{exo:b2:energy-geology:7}
$\lambda_5 - \lambda_8 = \qty{8.29725e-10}{yr^{-1}}$ en $r_0 = 0.7204/99.2742 = 0.0072567$.
(a) Op \qty{2.0}{\Ga} is $r = 0.0072567\,e^{1.6595} = 0.03814$, aandeel $0.03814/1.03814 =
\qty{3.67}{\percent}$; op \qty{1.8}{\Ga} is $r = 0.03231$, aandeel \qty{3.13}{\percent}. (b) $r =
3/97 = 0.030928$ en $t = \ln(0.030928/0.0072567)/(8.29725 \times 10^{-10}) = \qty{1.75}{\Ga}$.
(c) Het verarmde erts en normaal uranium zijn sindsdien met dezelfde factoren vervallen, zodat de verhouding van
hun aandelen $^{235}$U vandaag dezelfde is als vlak na de reacties: het erts bevatte toen de helft van het
normale aandeel, dus ongeveer de helft van zijn $^{235}$U was gespleten.
\end{solution}

\begin{solution}{exo:b2:energy-geology:8}
(a) $10^9/1.872 = \qty{5.34e8}{m^3}$; $10^9/621.4 = \qty{1.61e6}{m^3}$; $10^9/686.4 =
\qty{1.46e6}{m^3}$. (b) $7 \times 10^9/583 = \qty{1.20e7}{m^3}$ \ce{CO2}, die
$0.35 \times 0.30 = 0.105$ van het gesteentevolume vult: $1.20 \times 10^7/0.105 = \qty{1.14e8}{m^3}$ gesteente,
een gemiddelde dikte van $1.143 \times 10^8/(5 \times 10^6) = \qty{22.9}{m}$.
\end{solution}

\begin{solution}{exo:b2:energy-geology:9}
(a) $V = \qty{9e9}{m^3}$ en $q_R = 2.6 \times 10^6 \times 9 \times 10^9 \times 235 =
\qty{5.50e18}{J}$. (b) Over $30 \times \qty{3.156e7}{s} = \qty{9.47e8}{s}$: $R_g = 0.1$ geeft
\qty{581}{MW} en $0.25$ geeft \qty{1450}{MW} warmte. (c) $6 \times (10 \text{ tot } 20) = 60$ tot
\qty{120}{MW} elektriciteit, $60/581 = 0.10$ tot $120/581 = 0.21$ van het thermische vermogen: de meeste
warmte wordt afgevoerd, zoals in elke centrale die met een bescheiden temperatuurverschil werkt.
\end{solution}

\begin{solution}{exo:b2:energy-geology:10}
(a) De sterkte is $0.6 \times (80 - 39) = \qty{24.6}{MPa} > \qty{22}{MPa}$: stabiel. Bezwijken vraagt
$0.6\,\Delta p = 24.6 - 22$, $\Delta p = \qty{4.33}{MPa} = \qty{43}{bar}$, van dezelfde orde als de
\qty{32}{bar} die de aardbevingen van Denver begrensde. (b) $(24.6 - 24)/0.6 = \qty{1.0}{MPa} =
\qty{10}{bar}$: een breuk die al dicht bij bezwijken is, heeft maar een kleine stijging nodig. (c) Geïnjecteerde vloeistof verhoogt
de druk in een zone die zich door diffusie uitbreidt via doorlatend gesteente en breukzones, kilometers
ver over jaren; overal waar ze een breuk ontmoet die al dicht bij bezwijken is, volstaat een bescheiden stijging.
\end{solution}

\begin{solution}{exo:b2:energy-geology:11}
(a) Schrijf $x = e^{-u}$ met $u = (ART^2/\beta E)\,e^{-E/RT}$. Een omzettingsgraad van
\qty{10}{\percent} betekent $x = 0.9$, $u = \ln(10/9)$; \qty{90}{\percent} betekent $u = \ln 10$. Omdat
$\mathrm{d}\ln u/\mathrm{d}T = E/RT^2 + 2/T \approx E/RT^2$, is $\Delta(\ln u) \approx (E/RT^2)\,\Delta T$,
wat de formule geeft, met $\ln\ln 10 - \ln\ln(10/9) = 0.834 + 2.250 = 3.084$. (b) Met $T =
\qty{409.6}{K}$ is $RT^2/E = \qty{6.26}{K}$ en $\Delta T = \qty{19.3}{K}$; de exacte curve met één energie
geeft \qty{18.6}{K}, van 125.3 tot \qty{143.8}{\celsius}. (c) De doorgetrokken curve stijgt van
\qty{10}{\percent} rond \qty{122}{\celsius} tot \qty{90}{\percent} rond \qty{148}{\celsius}, over ongeveer
\qty{26}{K}: de spreiding van de activeringsenergieën verbreedt het omzettingsinterval.
\end{solution}

\begin{solution}{exo:b2:energy-geology:12}
(a) $2 \times 10^{10}/20 = \qty{1.0e9}{m^3}$ (\qty{1.0}{km^3}) en $2 \times 10^{10}/60 =
\qty{3.3e8}{m^3}$. (b) $2 \times 10^{10}/621.4 = \qty{3.22e7}{m^3}$. (c) $3.22 \times 10^7/0.02 =
\qty{1.61e9}{m^3}$ poriën, wat met $50 \times 0.30 = \qty{15}{m}$ poriënruimte per vierkante
meter $1.61 \times 10^9/15 = \qty{1.07e8}{m^2}$ vraagt, ongeveer \qty{107}{km^2}.
\end{solution}

\begin{solution}{pb:b2:energy-geology:1}
\noindent\textbf{Deel I.}
\textbf{1.} Op \qty{70}{\Ma}: $50 \times 60 = \qty{3000}{m}$, $15 + 25 \times 3.0 =
\qty{90}{\celsius}$. Vandaag: $3000 + 70 \times 30 = \qty{5100}{m}$, \qty{142.5}{\celsius}.
\textbf{2.} $60 \times 0.025 = \qty{1.5}{\celsius\per\Myr}$, daarna $30 \times 0.025 =
\qty{0.75}{\celsius\per\Myr}$.
\textbf{3.} 15--\qty{20}{\celsius}: \qty{3.33}{\Myr}; elk interval van 20 tot \qty{90}{\celsius}:
$10/1.5 = \qty{6.67}{\Myr}$; elk van 90 tot \qty{140}{\celsius}: $10/0.75 = \qty{13.3}{\Myr}$;
140--\qty{142.5}{\celsius}: \qty{3.33}{\Myr}.
\textbf{4.} Waarden per interval $\Delta t_n 2^n$, met lopende sommen: 15--20 ($n = -9$) 0.0065; 20--30
0.026 (0.033); 30--40 0.052 (0.085); 40--50 0.104 (0.189); 50--60 0.208 (0.397); 60--70 0.417
(0.814); 70--80 0.833 (1.65); 80--90 1.67 (3.31); 90--100 6.67 (9.98); 100--110 13.3 (23.3);
110--120 26.7 (50.0); 120--130 53.3 (103); 130--140 107 (210); 140--142.5 ($n = 4$) 53.3: TTI
vandaag 263.
\textbf{5.} 15 wordt gepasseerd in het interval 100--\qty{110}{\celsius}, bij $100 + 10 \times
(15 - 9.98)/13.33 = \qty{103.8}{\celsius}$, bereikt $(103.8 - 90)/0.75 = \qty{18.4}{\Myr}$ na
\qty{70}{\Ma}: op \qty{51.6}{\Ma} en $(103.8 - 15)/0.025 = \qty{3550}{m}$.
\textbf{6.} 75: \qty{124.7}{\celsius}, \qty{23.7}{\Ma}, \qty{4390}{m}. 160: \qty{135.3}{\celsius},
\qty{9.6}{\Ma}, \qty{4810}{m}.
\textbf{7.} Nee: met 263 is het voorbij het einde van de olievorming (160), sinds ongeveer \qty{10}{\Ma}.

\noindent\textbf{Deel II.}
\textbf{8.} $k = 9.21 \times 10^{13}\,e^{-26800/363.15} = \qty{8.20e-19}{s^{-1}}$. Het gesteente was
kouder dan \qty{90}{\celsius} gedurende de \qty{50}{\Myr} vóór \qty{70}{\Ma}, zodat de omgezette fractie
kleiner is dan $k\,t = 8.20 \times 10^{-19} \times 50 \times 3.156 \times 10^{13} = \num{1.29e-3}$.
\textbf{9.} $\beta = 0.75/(3.156 \times 10^{13}) = \qty{2.376e-14}{K/s}$, $T = \qty{415.65}{K}$:
$ART^2/(\beta E) = \num{2.498e28}$ en $e^{-E/RT} = \num{9.95e-29}$, dus $u = 2.486$, $x = 0.0833$ en
de omzettingsgraad is 0.917.
\textbf{10.} \qty{134.7}{\celsius}.
\textbf{11.} $70 - (134.7 - 90)/0.75 = \qty{10.4}{\Ma}$.
\textbf{12.} De TTI plaatst de maximale olievorming op \qty{23.7}{\Ma}, de kinetiek \qty{13}{\Myr} later,
omdat de verdubbelingsregel overeenkomt met $E \approx \qty{82}{kJ/mol}$ en de lage
temperaturen te zwaar weegt voor een kerogeen van \qty{223}{kJ/mol} (oefening~\ref{exo:b2:energy-geology:5}).
\textbf{13.} Eerder: een lagere energie verlaagt de temperatuur van halve omzetting.

\noindent\textbf{Deel III.}
\textbf{14.} Gesteente: $4 \times 10^8 \times 50 \times 2400 = \qty{4.8e13}{kg}$; potentieel $0.030
\times 4.8 \times 10^{13} = \qty{1.44e12}{kg}$; gevormd $0.92 \times 1.44 \times 10^{12} =
\qty{1.32e12}{kg}$.
\textbf{15.} $1.3248 \times 10^{12}/850 = \qty{1.56e9}{m^3}$.
\textbf{16.} $1.2 \times 10^7 \times 40 \times 0.22 \times 0.75/1.30 = \qty{6.09e7}{m^3}$.
\textbf{17.} $6.09 \times 10^7/1.56 \times 10^9 = 0.039$, \qty{3.9}{\percent}: de meeste gevormde olie
blijft in het brongesteente, gaat verloren langs de migratiewegen of ontsnapt naar het oppervlak.
\textbf{18.} $0.35 \times 6.09 \times 10^7 = \qty{2.13e7}{m^3}$.
\textbf{19.} Ongeveer twee derde van de olie blijft na primaire en secundaire winning in het reservoir;
injectie van \ce{CO2} zou nog eens 10 tot \qty{15}{\percent} van de aanwezige olie kunnen winnen.

\noindent\textbf{Deel IV.}
\textbf{20.} $15 + 25 \times 2 = \qty{65}{\celsius}$ en $0.101 + 1000 \times 9.807 \times 2000/10^6
= \qty{19.7}{MPa}$: beide boven het kritieke punt (\qty{31.0}{\celsius}, \qty{7.38}{MPa}), dus de
\ce{CO2} is superkritisch, met ongeveer \qty{686}{kg.m^{-3}}.
\textbf{21.} $2.13 \times 10^7 \times 1.30 = \qty{2.77e7}{m^3}$.
\textbf{22.} Poriënvolume $1.2 \times 10^7 \times 40 \times 0.22 = \qty{1.056e8}{m^3}$, en $E =
2.77 \times 10^7/1.056 \times 10^8 = 0.262$, ver boven de 1 tot \qty{4}{\percent} van zoute
formaties: de geproduceerde olie heeft een bekende, afgesloten ruimte vrijgemaakt.
\textbf{23.} $6 \times 10^8 \times 80 \times 0.20 \times 691 \times 0.02 = \qty{1.33e11}{kg}$,
\qty{133}{Mt}.
\textbf{24.} Met \qty{1}{Mt} per jaar ongeveer 133 jaar.
\textbf{25.} $M = \rho\,N_p\,B_o = 686 \times 2.77 \times 10^7 = \qty{1.90e10}{kg}$: het uitgeputte
reservoir kan \textbf{\qty[text-series-to-math=true]{19.0}{Mt}} \ce{CO2} opslaan.
\end{solution}
